\section{Layer III: Mathematical and Explanatory Structure}

\subsection{From Local Updates To Global Regimes}

The mathematical value of emergence language is not that it creates exact
closed-form formulas for everyday life. Its value is that it clarifies the
architecture of explanation.

\begin{proposition}[Micro-to-macro feedback]
If repeated local updates modify the cost of future transitions, then the
resulting macro-pattern is not merely descriptive. It becomes causally
relevant within the model because it changes what later trajectories are cheap,
stable, or accessible.
\end{proposition}

\begin{discussion}
This is the explanatory point behind causal-emergence language. Once a mature
regime changes later transition geometry, the regime deserves to be treated as
more than a summary label. It has become a control-relevant level of
description.
\end{discussion}

\subsection{Regimes, Basins, and Re-entry Cost}

The chapter does not require heavy equations, but it does require a stricter
mathematical picture. Let a behavioural regime be represented as a region of
state space with the following properties:
\begin{enumerate}[nosep]
  \item \textbf{entry cost:} how much coordination is required to enter it;
  \item \textbf{retention:} how likely trajectories are to remain inside it
        once entered;
  \item \textbf{re-entry economy:} how cheaply the system can return after
        interruption;
  \item \textbf{competitive suppression:} how strongly rival routes are made
        less attractive by the established regime.
\end{enumerate}

Framework formation can then be interpreted as a process that progressively:
\begin{itemize}[nosep]
  \item lowers entry cost for one route,
  \item deepens or sharpens its basin of attraction,
  \item stores route memory through repeated coordinated use,
  \item increases the interference cost of competing paths.
\end{itemize}

\subsection{Why Emergence Is Not Just Aggregation}

Simple aggregation would mean that many repetitions merely add up. Emergence in
the present chapter is stronger than that. The repeated acts do not just
accumulate; they reorganize the future problem.

\begin{lemma}[Reorganization criterion]
If repeated local events change the future transition landscape itself, then
the resulting higher-level pattern is doing more than summarizing the events.
\end{lemma}

\begin{proof}
Suppose the local events merely add up without reorganizing later choice. Then
the future decision problem is structurally unchanged, and the repeated events
have not created a new explanatory level. But if the events lower one route's
re-entry cost, strengthen its retention, or increase interference against
rival routes, then the future problem is different from what it would have been
without the history. The macro-pattern has therefore become model-relevant.
\end{proof}

\subsection{Mathematical Upgrade Path}

This chapter is intentionally lighter than Chapter~\ref{ch:advanced-dynamics},
but it prepares that chapter in three ways:
\begin{enumerate}[nosep]
  \item metastability clarifies why useful frameworks must persist without
        becoming immutable;
  \item attractor language clarifies how repeated use can turn one route into a
        default basin;
  \item operator and control viewpoints later help describe how interventions
        act on a history-sensitive system rather than a blank slate.
\end{enumerate}
